Predicting the resonant modes of three-dimensional radio-frequency (RF) cavities---their frequencies, electric
fields and figures of merit such as $\Qz$, $\RoQ$ and the peak surface fields---normally requires one
finite-element eigensolve per geometry. We present a neural operator that does not regress these quantities
directly. It maps the geometry, represented by the quadrature points of a curved tetrahedral mesh, to a learned
basis of vector fields; a fixed linear map turns the basis into degrees of freedom of a second-order N\'ed\'elec
space, a kernel projection removes the discrete gradients, and a Rayleigh--Ritz step returns the modes. The
predicted eigenvalues are therefore upper bounds of the discrete eigenvalues of whichever conforming space the
model runs in, and no spurious zero modes arise. Single-sample fits reach relative eigenvalue errors of
$10^{-6}$--$9\times10^{-6}$. % src: docs/28 §3
We build an adaptive third-order label generator on curved meshes, validated against analytic pillbox modes
(relative frequency error $6\times10^{-7}$ at 157k degrees of freedom), % src: docs/31 §1
and show that low-fidelity labels---one second-order solve on the same model mesh---cost $4.3\times$ less
while differing from the adaptive labels by at most $0.047\,\%$ in frequency and $0.32\,\%$ in field. % src: docs/34 §8.1
In a controlled small-scale study on one cavity family, models trained on low-fidelity labels were at least as
accurate on held-out high-fidelity geometries as models trained on high-fidelity labels of equal or higher
cost, and their errors against low- and high-fidelity labels differ by $0.02$ percentage points. % src: docs/34 §8.2
